\section{Related Work}
\label{sec:related}

The PolieBotics Truth Beam draws on several established lines of work, but combines them in a way that none of them individually addresses. We organize the discussion around five threads: active illumination and structured light; media provenance and cryptographic content signing; diffusion models as anomaly and likelihood detectors; conditioning mechanisms for image-to-image networks; and verifiable distributed randomness. In each case we delineate what we inherit from prior art and what is specific to a projector--camera capture chain whose authenticity rests on a \emph{physical} coupling between an emitted pattern and the captured scene.

\subsection{Structured light and active illumination}
The capture front end is, mechanically, a structured-light system: a projector emits a known pattern and a synchronized camera records the scene under that illumination. This places the apparatus in the same hardware family as the projector--camera systems of the underlying patent \cite{patent2025}. Classical structured light uses the projected pattern to recover geometry (depth, surface normals) by triangulating known stripes or codes, and the pattern is chosen for ease of decoding~\cite{geng2011}. Our use of active illumination is different in intent. The emitted tiles $E_t$ are not designed to be geometrically decodable; they are deterministic fractional Brownian motion (fBm) fields derived per-channel from a cryptographic state via BLAKE3-XOF expansion (\num{43110} bytes per channel decomposed into four octaves), rendered as \SI{8}{bit} RGB and projected at $1920\times1080$, then captured as an \SI{8}{bit} BayerRG8 raw frame $C_t$ of exactly \num{24472000} bytes. The pattern's role is to act as a per-frame, content-unpredictable optical probe: because $E_t$ is fixed only once the chain state $S_t$ is known, and because $S_{t+1}$ folds in a hash of the raw capture $C_t$, the illumination at each step is bound to the recorded history of the session. The verification question is therefore not ``what is the geometry?'' but ``is this capture physically consistent with the illumination that the chain says was projected?''. Active illumination has also been used directly for video authenticity: Gerstner and Farid detect real-time deep fakes in video calls by analysing the physical response to time-varying screen illumination~\cite{gerstner2022active}, and GOTCHA~\cite{mittal2024gotcha} formalises a challenge--response protocol that exploits the live, interactive setting against real-time deepfake pipelines. Those systems issue face-centric, screen-scale challenges chosen for a human-in-a-call scenario; the Truth Beam generalises the stance to a full-scene structured-light challenge stream in which every probe is derived from -- and cryptographically committed to -- the recording's own hash-chained history, and in which the response (the raw capture) feeds back to determine the next challenge. The protocol is \SI{8}{bit} end-to-end throughout this loop; the sensor is a Sony IMX540 whose \SI{12}{bit} ADC is configured to deliver \SI{8}{bit} output, so the structured-light substrate we exploit lives entirely in the \SI{8}{bit} optical regime.

\subsection{Media provenance and cryptographic content signing}
A large and growing body of work secures media authenticity by attaching cryptographic signatures or provenance manifests to a file after capture, as in C2PA-style content credentials~\cite{c2pa} and trusted-capture pipelines. These schemes answer the question ``was this bitstream produced by a key-holding device and not altered since?''. They are powerful, but the trust they confer is fundamentally a statement about \emph{bits and keys}, not about the \emph{physical event} that the bits purport to depict: a signing key that is exfiltrated, or a capture device fed a synthetically rendered frame, will sign a forgery exactly as readily as a genuine capture. A complementary passive line identifies the \emph{source camera} from its sensor's photo-response non-uniformity (PRNU) fingerprint~\cite{lukas2006prnu}; that attests which device produced a frame, not what physical scene interaction produced it. The Truth Beam takes a complementary approach to both. Rather than (only) signing the output or fingerprinting the device, it makes the recording medium itself tamper-evident by closing an optical feedback loop. At each step the 32-byte chain state $S_t$ is BLAKE3-XOF-expanded \cite{blake3} into the emission $E_t$; $E_t$ is projected and captured as $C_t$; the raw Bayer bytes are hashed (\texttt{bayer\_blake3}); and that hash is folded back into $S_{t+1}$ together with a 28-byte capture-meta struct and the beacon fields. Substituting a forged capture for $C_t$ yields a different Bayer hash, hence a different $S_{t+1}$, hence a different emission $E_{t+1}$ --- a discrepancy any verifier holding a sample of the chain can detect. The two chain variants used here, v9 (session D2) and v10 (session V10, which additionally commits a 32-byte \texttt{ai\_payload\_root} for AI-agent response frames), differ in their domain-separated transition formula but share this closed-loop property. The distinguishing claim relative to signature-only provenance is thus not cryptographic strength per se, but that authenticity is anchored to a learnable physical coupling between $C_t$ and $E_t$ that a downstream forger must reproduce optically rather than merely sign. We are deliberately conservative about how far this extends: the present evidence establishes the coupling on a single rig across two sessions and two chain protocols, and does not address signature-infrastructure threats such as key compromise.

\subsection{Diffusion models as anomaly and likelihood detectors}
Our verifier is built on a denoising diffusion probabilistic model. We adopt the DDPM noise-prediction formulation of Ho et al.\ \cite{ho2020ddpm} and the cosine noise schedule of Nichol and Dhariwal \cite{nicholdhariwal2021} (offset $s=0.008$, $T=1000$). A key departure from the generative use of such models is that the verifier never samples or denoises: it is run purely as a scoring function, computing the noise-prediction residual $\lVert \hat{\epsilon}-\epsilon\rVert^2$ at five fixed timesteps $(50,150,300,500,750)$ with $K_{\text{NOISE}}=4$ noise draws per $(C,E,t)$. Using the reconstruction or likelihood behavior of a generative model as an anomaly score is a recognized strategy --- diffusion priors have recently been applied to general image-forgery detection and localization~\cite{yu2024diffforensics}, where the model supplies a learned notion of image typicality; our verifier differs in being conditioned on a per-frame physical challenge rather than scoring an image in isolation --- but it carries a well-documented hazard: deep generative models can assign \emph{higher} likelihood (lower reconstruction error) to out-of-distribution inputs than to in-distribution data, the ``overspread'' phenomenon characterized by Nalisnick et al.\ \cite{nalisnick2019}. This caveat is not incidental to our work; it is the mechanism most consistent with one of our most important negative findings. Under proper Mask R-CNN \cite{he2017maskrcnn} body segmentation, scoring \emph{body-only} content with the $\epsilon$-MSE residual produces AUROC values that collapse to chance or below (E2$_{\mathrm{seg}}$ AUROC $0.009$--$0.538$ across conditions; values below $0.5$ indicate that, restricted to the body region, F-A v1 fake content sits nearer the verifier's conditional-likelihood mode than genuine training-distribution samples). We frame this strictly in the distribution-distance / likelihood-overspread language of \cite{nalisnick2019} rather than as any claim that ``fakes look more real than real''. Crucially, the discriminative signal is not absent: it lives off-body (E3$_{\mathrm{seg}}$ AUROC $=1.000$ across all twelve session$\times$condition cells) and, on the body region, survives at low spatial frequencies (a Gaussian-blurred residual at $\sigma=4$ retains body-only AUROC $\approx 0.998$) where the high-frequency overspread that dominates the raw $\epsilon$-MSE no longer masks it. This work thus sits squarely in the lineage of likelihood-based OOD detection while taking the overspread caveat as a first-class design constraint rather than a footnote.

\subsection{Conditioning the verifier on the emission}
The verifier must score a capture \emph{relative to} the specific emission the chain prescribes, so the architecture needs a mechanism to inject $E_t$ as a conditioning signal. We adopt the spatial, additive conditioning pattern of ControlNet \cite{zhang2023controlnet}: rather than concatenating $E_t$ to the camera input, the emission is passed through a four-stage strided HintEncoder ($/16$ downsampling) and \emph{added} into the U-Net's down-path features at each scale through zero-initialized $1\times1$ ``zero-conv'' adapters, so that at training step zero the network is unconditional and the emission's contribution ramps up smoothly. This is preferable to coarser global modulation schemes such as FiLM \cite{perez2018film}, whose channel-wise affine transform cannot carry the spatially localized emission structure that the substrate verification depends on. The choice is consequential: an architecture- and data-matched \emph{shuffled} control sits at chance (AUROC $\approx 0.5006$, $\Delta_{\text{wrong}}\approx 10^{-7}$), confirming that the perfect separation observed in the main run reflects genuine $C$--$E$ coupling learned through the conditioning path rather than a generic $E$-statistic or an architecture artifact. The same ControlNet-style spatial-hint design recurs in our red-team attacker (F-A v1, a $42.3$M-parameter EditorControlNet over a ConvNeXt-Tiny encoder), so that conditioning mechanism appears on both sides of the verifier/attacker dialogue.

\subsection{Verifiable distributed randomness}
Finally, the chain incorporates an external source of public, verifiable randomness. At each transition the writer folds in a drand quicknet beacon \cite{drandbeacon}: an 8-byte big-endian round number and the 32-byte round randomness (the SHA-256 of that round's BLS signature). Including a publicly auditable beacon timestamps the chain against an independent randomness source that the recorder cannot have predicted in advance, strengthening the temporal-ordering guarantee beyond what a purely self-referential hash chain provides. We are explicit about the boundary of this contribution: the chain transition itself does not re-verify the beacon signature (the writer is required to verify it locally before folding it in), and when the beacon is unavailable a sentinel (round number $0$, zero randomness) is folded in verbatim, which the verifier notes as a \texttt{drand\_availability} status but does not use to gate a PASS/FAIL decision. The beacon thus contributes auditability and unpredictability to the provenance record rather than serving as a hard cryptographic gate.

\subsection{Positioning}
In summary, the Truth Beam is not a new structured-light geometry method, not a content-signing scheme, and not a generative diffusion model: it borrows the projector--camera apparatus from the first, the tamper-evidence goal from the second, the scoring backbone from the third, and the conditioning mechanism and randomness anchor from the fourth and fifth threads. Its distinctive contribution is to bind a cryptographic hash chain to a physical optical coupling and to learn that coupling with a diffusion-residual verifier, while treating the OOD-overspread failure mode of likelihood-based detection \cite{nalisnick2019} as a central, empirically confirmed phenomenon rather than an afterthought. In grounding authenticity in a hard-to-reproduce physical interaction rather than solely in keys and signatures, the approach is also a descendant of physical one-way functions and optical PUFs~\cite{pappu2002}, a line that has since extended to sensor-level PUFs for camera authentication~\cite{zheng2020edpuf} --- with the distinction that here the challenge--response pair is generated by the recording's own hash-chained history and is read out by a learned verifier over a live scene rather than by a fixed analytic comparison over a static token or device fingerprint. We restate the scope that governs every comparison above: all results are same-rig, cross-session (and same-rig cross-protocol) only, established against the trained F-A v1 attacker; cross-rig, cross-camera, cross-projector, and cross-subject generalization, and robustness against the stronger (untrained) F-A v2 attacker, remain future work.
